% To demonstrate the necessity of our proposed approach and
% highlight the intrinsic limitations of existing tools, this
% section presents three illustrative examples. Specifically,
% we examine three representative scenarios that pose significant
% challenges to current verification techniques: \circled{1} proofs
% involving parameterized bitwidths, \circled{2} complex cases where
% existing tools suffer from timeout failures, and \circled{3}
% verification tasks related to loop structures. Through these
% examples, we elucidate how \proj effectively addresses
% the shortcomings of \alivetwo.

This subsection presents two concrete verification tasks from our benchmark
suite: one valid transformation and one invalid transformation. For each
task, we describe the source and target fragments, the resulting proof
obligation, and the certified outcome produced by \proj.

\subsubsection{Certifying a Valid Transformation}

The first example is extracted from an \llvm missed optimization
issue~\cite{illustrative_example1_issue}: an
arbitrary-bitwidth generalization of a masked-bit-test rewrite
that \alivetwo cannot verify.
In the source fragment, the \texttt{nsw} flags guarantee that
\texttt{p2 = 2 * a * b} is a well-defined even integer. Therefore
\texttt{p2 | 1} only sets the low bit and preserves the sign bit,
making \texttt{icmp slt (p2 | 1), 1} equivalent to testing
\texttt{p2 < 0}. Since \texttt{p2 = 2 * a * b} without signed
overflow, this is equivalent to \texttt{a * b < 0}. The source then
returns \texttt{a * b} when the condition holds and \texttt{a}
otherwise, matching the target's \texttt{select}.
The notation \texttt{inputs a : w, b : w} denotes transformation inputs
where \texttt{a} and \texttt{b} are integers represented as bit-vectors
of width \texttt{w}.

\begin{lstlisting}[
  basicstyle=\ttfamily\small,
  xleftmargin=1em,
  aboveskip=0.5\baselineskip,
  belowskip=0.5\baselineskip,
  morekeywords={Source,Target,assume,return,inputs}
]
Source:
  inputs a : w, b : w
  a2     : w  = shl nsw a, 1
  p2     : w  = mul nsw b, a2
  p2_or1 : w  = or p2, 1
  cond   : i1 = icmp slt p2_or1, 1
  factor : w  = select cond, b, 1
  r      : w  = mul a, factor
  return r

Target:
  inputs a : w, b : w
  p    : w  = mul b, a
  cond : i1 = icmp slt p, 0
  r    : w  = select cond, p, a
  return r
\end{lstlisting}

\proj certifies this transformation through its proof pipeline.
These two fragments are first translated into a refinement theorem
universally quantified over the bitwidth:
\[
\forall (\texttt{w1 : Nat}),\;
  \forall (\texttt{h : 0 < w1}),\;
  (\texttt{Source w1 h}) \sqsubseteq (\texttt{Target w1 h})
\]
In this theorem, \texttt{w1} denotes the symbolic bitwidth
in the source and target fragments, and \texttt{h} is
the proof parameter enforcing \(0 < \texttt{w1}\).

\proj then decomposes the refinement goal into smaller
obligations over bit-vectors and LLVM definedness conditions.
To avoid redundant proof efforts, the pipeline discharges these
obligations by matching recurring subgoals against a reusable
lemma library and instantiating the lemmas with current
operands and bitwidths. In this example, the refinement branch
applies a lemma that interprets a no-overflow left shift
as a signed multiplication by two.
For subgoals absent from the lemma library and any remaining
proof steps, the \llm synthesizes the proofs from scratch.
Moreover, the proof separately accounts for undefined behavior,
such as poison value inputs and poison values generated by
the \texttt{nsw}-annotated operations. The final refinement
lifts the equality of returned integer values to the
theorem-level source-target refinement. Once the proof succeeds,
\proj extracts reusable lemmas from it and adds them to the library
for future proofs. It then terminates the counterexample branch for this
transformation.

\subsubsection{Refuting an Invalid Transformation}

The second example comes from Hydra~\cite{hydra}, which generalizes a
real-world \llvm transformation~\cite{illustrative_example2_issue}
over symbolic bit-widths \texttt{w1} and
\texttt{w2}, under the precondition $0 < \texttt{w1} < \texttt{w2}$
and \texttt{C2 = w1 - 1}.
In the source fragment, the precondition makes \texttt{C2} the
sign-bit index of \texttt{x}. Thus, \texttt{lshr x, C2} extracts
the sign bit to construct an all-ones or all-zero mask,
conditionally yielding \texttt{C1 - sext(x)} for negative
\texttt{x} and zero otherwise. The target attempts to replace
this mask construction into a single signed comparison against
a shifted threshold.
This transformation uses two symbolic bitwidths: \texttt{x} and
\texttt{C2} have width \texttt{w1}, while \texttt{C1} has width
\texttt{w2}.

\begin{lstlisting}[
  basicstyle=\ttfamily\small,
  xleftmargin=1em,
  aboveskip=0.5\baselineskip,
  belowskip=0.5\baselineskip,
  morekeywords={Source,Target,assume,return,inputs}
]
Source:
  inputs x : w1, C1 : w2, C2 : w1
  assume C2 == w1 - 1
  x_lshr      : w1 = lshr x, C2
  x_lshr_zext : w2 = zext x_lshr from w1 to w2
  neg_mask    : w2 = sub 0, x_lshr_zext
  x_sext      : w2 = sext x from w1 to w2
  sub         : w2 = sub C1, x_sext
  result      : w2 = and neg_mask, sub
  return result

Target:
  inputs x : w1, C1 : w2, C2 : w1
  assume C2 == w1 - 1
  x_sext    : w2 = sext x from w1 to w2
  sub       : w2 = sub C1, x_sext
  C2_plus_1 : w1 = add C2, 1
  threshold : w1 = shl 1, C2_plus_1
  cond      : i1 = icmp slt x, threshold
  result    : w2 = select cond, sub, 0
  return result
\end{lstlisting}

\proj refutes this transformation through its counterexample pipeline.
These two fragments are first encoded as parameterized Lean
definitions. The counterexample branch leverages \llm to propose a
counterexample candidate with concrete widths and inputs satisfying
the precondition. Then the counterexample candidate is used to prove the
existential non-refinement theorem:
\[
\begin{aligned}
&\exists (\texttt{w1 w2 : Nat}),\;
  \exists (\texttt{h : w1 < w2}),\;
  \exists (\texttt{h1 : 0 < w1}),\\
&\qquad
  \neg((\texttt{Source w1 w2 h h1}) \sqsubseteq
       (\texttt{Target w1 w2 h h1}))
\end{aligned}
\]
Here, \texttt{w1} and \texttt{w2} are symbolic bitwidths, while
\texttt{h} and \texttt{h1} enforce \(\texttt{w1} < \texttt{w2}\)
and \(0 < \texttt{w1}\), respectively.

In this example,
\proj proposes a counterexample candidate \texttt{w1 = 2},
\texttt{w2 = 3}, with \texttt{x = 2}, \texttt{C1 = 0}, and \texttt{C2 = 1}.
This candidate satisfies the precondition because $0 < 2 < 3$ and
\texttt{C2 = w1 - 1}. On the source side, the fragment returns
the concrete value \texttt{2}. On the target side, however,
\texttt{C2\_plus\_1} is \texttt{2} at bitwidth two, so
\texttt{shl 1, 2} shifts by an amount equal to the bitwidth
and produces poison value under the \llvm semantics.
Since a concrete source result cannot be refined by a poison
target result, this valuation violates \(\texttt{Source} \sqsubseteq \texttt{Target}\)
and therefore constitutes a counterexample. \proj certifies this
counterexample by lifting the concrete source-defined/target-poison
discrepancy to the \llvm refinement relation, thereby establishing
the fixed-width non-refinement instance for \texttt{w1 = 2} and
\texttt{w2 = 3}. This instance is then used to discharge the
existential non-refinement theorem.
Upon successful refutation, \proj reports the counterexample
and terminates the refinement branch for this transformation.
